\documentclass[12pt]{article}
\usepackage[T2A]{fontenc}
\usepackage[utf8]{inputenc}
\usepackage[russian]{babel}
\usepackage{amsmath, amssymb, amsfonts}
\usepackage{geometry}
\usepackage{fancyhdr}
\usepackage{microtype}
\usepackage{titlesec} % Пакет для управления стилем заголовков
\usepackage{enumitem}
\usepackage{tikz}
\usepackage{tikz-3dplot}
\usetikzlibrary{calc, intersections, through, backgrounds, arrows.meta}

% Настройка Section: по центру, жирный, крупный
\pagestyle{fancy}
% Настройка полей
\geometry{a5paper, margin=1cm, top=2cm, headheight=15pt}

% --- ЖЕСТКАЯ КОРРЕКЦИЯ ЦЕНТРИРОВАНИЯ ---
\pagestyle{fancy}
\fancyhf{} % Полная очистка всех полей (убирает невидимые отступы)
\fancyhead[C]{\thepage} % Номер строго по центру
\renewcommand{\headrulewidth}{0pt} % Убираем линию
\fancyfoot{} % Низ пустой

% Убираем внутренние отступы колонтитулов, которые могут давать перекос
\setlength{\headwidth}{\textwidth} 

% Переопределяем системный стиль plain, чтобы он не сбивал настройки
\makeatletter
\let\ps@plain\ps@fancy
\makeatother

\parindent=0pt
\parskip=0.6em
\titleformat{\section}
  {\normalfont\bfseries\centering}{}{0pt}{}

% Настройка Subsection: отступ 1см, жирный, обычный размер
% {команда}[формат]{стиль}{метка}{отступ_от_метки}{перед_заголовком}[после_заголовка]
\titleformat{\subsection}
  {\normalfont\bfseries}{\hspace{1cm}}{1pt}{}
\titlespacing*{\subsection}{1cm}{2ex plus 1ex minus .2ex}{0.5ex plus .2ex}

% Убираем автоматическую нумерацию, чтобы заголовки были как в оригинале
\setcounter{secnumdepth}{0}
\AddEnumerateCounter{\asbuk}{\russian@asbuk@alph}{щ} % section 8.1
\setlist{wide, nosep, font=\rmfamily\normalfont}
\setlist[enumerate, 2]{label=\asbuk*)}
\setlist[enumerate, 3]{label=\asbuk*)}

\begin{document}
	\section{Глава 5. Разное}
	\subsection{5.1. Геометрия и аналитика}
	В курсе неоднократно проводилась мысль о том, что геометрические задачи можно двояко -- геометрически и аналитически. Какие решения проще? Этот спор невозможно разрешить. В нижеследующем цикле предлагается собрание классических (и не только классических) задач и каждый для себя может устроить состязание, достигая цели наиболее простым путем, и вынести приговор о том, какой из путей -- аналитический или геометрический -- быстрее ведет к цели.
	\begin{enumerate}
		\item Докажите, что центр тяжести, центр описанного круга и ортоцентр треугольника лежат на одной прямой (прямая Эйлера).
		\item Пусть $ABCD$ -- произвольный четырехугольник, $E$ -- точка пересечения $AB$ и $CD$, $F$ -- точка пересечения $BC$ и $AD$, $M$, $N$, $O$ -- середины $AC$, $BD$ и $EF$ соответсвенно. Доказать, что $M, N, O$ лежат на одной прямой (прямая Гаусса).
		\item Дан треугольник $ABC$. На его сторонах $AB, BC$ и $CA$ взяты точки $C_1, A_1$ и $B_1$. Доказать, что прямые $AA_1, BB_1$ и $CC_1$ пересекаются в одной точке тогда и только тогда, когда 
		\[
			\frac{|BA_1|}{|CA_1|} \cdot \frac{|CB_1|}{|AB_1|} \cdot \frac{|AC_1|}{|BC_1|} = 1 \quad \text{(теорема Чевы)}.
		\]

		\item Дан треугольник $ABC$. На сторонах $AB$ и $BC$ взяты точки $C_1$ и $A_1$, а точка $B_1$ взята на продолжении стороны $CA$. Доказать, что точки $A_1, B_1$ и $C_1$ лежат на одной прямой тогда и только тогда, когда 
		\[ 
			\frac{|BA_1|}{|CA_1|} \cdot \frac{|CB_1|}{|AB_1|} \cdot \frac{|AC_1|}{|BC_1|} = 1 \quad \text{(теорема Менелая)}.
		\]

		\item Дан четырехугольник $ABCD$, $E$ — точка пересечения $CD$ и $AB$, $F$ — точка пересечения $BC$ и $AD$. Доказать, что окружности, описанные около $AED, ABF, BCE$ и $CDF$ пересекаются в одной точке.

		\item Пространственный неплоский шестиугольник имеет три пары параллельных противоположных сторон. Докажите, что в каждой из этих пар стороны равны (из вступительного экзамена в Математический Колледж Независимого Московского Университета, 1992).

		\item Провести через точку $P$ внутри данного угла отрезок $MN$ с концами на сторонах угла так, чтобы произведение $|MP| \cdot |PN|$ было наименьшее (из вступительного экзамена в МК НМУ, 1992).

		\item Дан выпуклый четырехугольник $ABMC$, в котором $|AB| = |BC|$, $\angle BAM = 90^\circ$, $\angle ACM = 150^\circ$. Докажите, что $AM$ — биссектриса $\angle BMC$ (задача 56-й московской олимпиады).

		\item К стороне $AB$ треугольника $ABC$ внешним образом построен квадрат с центром $O$. Точки $M$ и $N$ — середины сторон $AC$ и $BC$ соответственно, а длины сторон равны $a$ и $b$. Найти максимум $|OM| + |ON|$, когда угол $ACB$ меняется (задача 56-й московской олимпиады).

		\item Дан треугольник $ABC$, $r$ — радиус вписанной, $R$ — радиус описанной окружности. Доказать, что $2r \le R$ (Эйлер).

		\item Дан единичный круг. Через заданную точку $F$ на диаметре $AB$ провести хорду $CD$ так, чтобы площадь четырехугольника $ACBD$ была максимальной.

		\item Дан угол $BAC$ и две точки $M$ и $N$ внутри него. Провести отрезок $DE$ через $M$ (с помощью циркуля и линейки), чтобы площадь четырехугольника $ADNE$ была минимальной.

		\item Среди пирамид с данными основанием и высотой найти пирамиду с наименьшей боковой поверхностью.
	\end{enumerate}

	\subsection{5.2. Геометрия и алгебра}
	\begin{enumerate}
		\item Сколько элементов в группе \textit{вращений} правильного $n$-угольника и сколько в группе его \textit{изометрий}?

		\item Образует ли группу множество всех \textit{переворотов} правильного $n$-угольника?

		\item Доказать, что группа вращений правильного $n$-угольника изоморфна группе вычетов по модулю $n$.

		\item Доказать, что при $n > 2$ группа изометрий правильного $n$-угольника не коммутативна.

		\item Сколько разных подгрупп содержат группы вычетов по модулям: $3, 4, 5, 6, 7, 8, 9, 10, 12, 15, 120$? Каковы их порядки?

		\item Сколько разных подгрупп содержит группа изометрий правильного $n$-угольника при $n = 3, 4, 5, 6, 7, 8, 15, 120$? Каковы их порядки?

		\item Доказать, что любая группа, состоящая из простого числа элементов, изоморфна группе вычетов по простому модулю.

		\item Доказать, что группа целых чисел не изоморфна группе целых векторов на плоскости.

		\item Сколько элементов содержат группы вращений пяти правильных многогранников? Какие из этих групп изоморфны между собой?

		\item Доказать, что все группы вращений правильных многогранников не коммутативны.
		
		\item Доказать, что при $n > 2$ группа перестановок длины $n$ не коммутативна. Сколько элементов в этой группе?

		\item Сколько разных подгрупп содержит группа перестановок длины $n$?

		\item Доказать, что все \textit{четные} перестановки длины $n$ образуют подгруппу в группе перестановок. Сколько элементов в этой подгруппе?

		\item Доказать, что каждая из групп вращений правильных многогранников изоморфна одной из групп перестановок (всех -- данной длины -- или только четных). Какой геометрический смысл можно придать переставляемым символам в этих трех случаях?

		\item Доказать, что группа $\mathbb{Z}_{6}$ изоморфна прямой сумме групп $\mathbb{Z}_2 \oplus \mathbb{Z}_3$, а группа $\mathbb{Z}_{4}$ не изоморфна сумме $\mathbb{Z}_2 \oplus \mathbb{Z}_2$. 
		При каких условиях возможен изоморфизм $\mathbb{Z}_k \oplus \mathbb{Z}_n \cong \mathbb{Z}_{kn}$?

		\item Сколько существует попарно не изоморфных коммутативных групп из 8 элементов?

		\item Доказать, что группа кватернионных единиц не изоморфна группе изометрий квадрата.

		\item Доказать, что всякая конечная группа вкладывается в группу перестановок. Какой длины перестановки нужно брать для этого?

		\item Построить пример бесконечной группы, все элементы которой имеют порядок 2.

		\item Доказать, что свободные коммутативные группы с разным числом образующих не изоморфны между собой.

		\item Доказать, что не изоморфны группы $\mathbb{Z}_2 \oplus \mathbb{Z}_n$ и $\mathbb{Z}_2 \oplus \mathbb{Z}_m$, если $m \ne n$.

		\item Для произвольного множества $A$ ввести в множестве всех его подмножеств такую операцию, относительно которой оно станет коммутативной группой, причем порядок каждого элемента будет равен 2.

		\item Сколько элементов порядков $2, 3, 4, 5, 6, 8$ содержится в группе вращений окружности $SO_2$?

		\item Приведите пример элемента бесконечного порядка в группе $SO_2$.

		\item Есть ли в $SO_2$ свободная коммутативная подгруппа с 10 образующими?

		\item Сколько элементов порядка 2 содержится в группе $O_2$? Сколько в ней элементов порядка $3, 4, 5, 8$?

		\item Есть ли в $O_2$ подгруппы, изоморфные $\mathbb{Z}_2 \oplus \mathbb{Z}_2$? $\mathbb{Z}_2 \oplus \mathbb{Z}_2 \oplus \mathbb{Z}_2$?

		\item Найти группы симметрий правильного треугольника, тетраэдра, куба, их подгруппы и нормальные делители.
	\end{enumerate}
	\subsection{5.3. Геометрия и топология}
	\begin{enumerate}

		\item Как формулируется теорема Эйлера о многогранниках в $n$-мерном пространстве?

		\item Введите естественную метрику во множестве всех прямых на плоскости и докажите, что это пространство гомеоморфно листу Мёбиуса.

		\item Докажите эквивалентность следующих определений проективной плоскости:
		\begin{enumerate}
			\item двумерный диск со склеенными парами диаметрально противоположных точек;
			\item двумерная сфера со склеенными концами всех диаметров;
			\item множество всех прямых в $\mathbb{R}^3$, проходящих через начало.
		\end{enumerate}

		\item Вложить проективную плоскость в $\mathbb{R}^4$ без самопересечений.

		\item Разбить проективную плоскость на шесть попарно граничащих областей, гомеоморфных диску.

		\item Разбить шар на семь попарно граничащих областей, гомеоморфных диску.

		\item Доказать, что дополнение к кругу в проективной плоскости есть лист Мёбиуса.

		\item Доказать, что трехмерное проективное пространство гомеоморфно группе вращений трехмерного евклидова пространства.

		\item Доказать, что однополостный гиперболоид в проективном пространстве диффеоморфен тору.

		\item Доказать, что квадрика в $\mathbb{C}P^3$ диффеоморфна $S^2 \times S^2$.

		\item Доказать, что пространство окружностей в пространстве эллипсов имеет коразмерность 2.
	\end{enumerate}
\end{document}
